\documentclass[11pt,a4paper]{article}

% ── Packages ────────────────────────────────────────────────────────
\usepackage[UTF8]{ctex}           % Chinese support
\usepackage{amsmath,amssymb,amsthm}
\usepackage{mathtools}
\usepackage{tikz-cd}              % Commutative diagrams
\usepackage{xcolor}
\usepackage{hyperref}
\usepackage{geometry}
\usepackage{enumitem}
\usepackage{booktabs}
\usepackage{fontspec}

\geometry{margin=2.5cm}
\hypersetup{
    colorlinks=true,
    linkcolor=blue!60!black,
    citecolor=blue!40!black,
    urlcolor=blue!60!black,
}

% ── Typography ──────────────────────────────────────────────────────
\setlength{\emergencystretch}{3em}   % Extra tolerance for overfull boxes
\setlength{\hfuzz}{2pt}              % Ignore trivial overflows
\tolerance=2000                       % Allow significant stretching for math+Chinese
\raggedbottom                         % Avoid underfull vbox
\allowdisplaybreaks                   % Allow multi-line equations to break
\binoppenalty=10000                   % Discourage inline binop breaks
\relpenalty=10000                     % Discourage inline relation breaks

% ── Theorem environments ────────────────────────────────────────────
\theoremstyle{definition}
\newtheorem{theorem}{定理}[section]
\newtheorem{lemma}[theorem]{引理}
\newtheorem{proposition}[theorem]{命题}
\newtheorem{corollary}[theorem]{推论}
\theoremstyle{definition}
\newtheorem{definition}[theorem]{定义}
\newtheorem{axiom}[theorem]{公理}
\theoremstyle{remark}
\newtheorem{remark}[theorem]{注释}
\newtheorem{example}[theorem]{例}

% ── Category theory macros ──────────────────────────────────────────
\newcommand{\cat}[1]{\mathbf{#1}}
\newcommand{\catWFlow}{\cat{WFlow}}
\newcommand{\catCap}{\cat{Cap}}
\newcommand{\catDet}{\cat{Det}}
\newcommand{\catNonDet}{\cat{NonDet}}
\newcommand{\catComp}{\cat{Comp}}
\newcommand{\catBool}{\cat{Bool}}
\newcommand{\catSet}{\cat{Set}}
\newcommand{\id}{\mathrm{id}}
\newcommand{\ob}[1]{\mathrm{Ob}(#1)}
\newcommand{\homset}[3]{\mathrm{Hom}_{#1}(#2, #3)}
\newcommand{\functor}[3]{#1: #2 \to #3}
\newcommand{\natarrow}{\Rightarrow}
\newcommand{\adjunction}{\dashv}
\newcommand{\tensor}{\otimes}
\newcommand{\monadT}{T}
\newcommand{\pairMonad}{P}
\newcommand{\Yoneda}{\mathcal{Y}}
\newcommand{\pairCoupling}{\bowtie}     % Pair coupling (Harness+LLM nesting)
% Semantic brackets for denotation
\newcommand{\llbracket}{[\![}
\newcommand{\rrbracket}{]\!]}

% ── Title ───────────────────────────────────────────────────────────
\title{
    {\Huge The Pair 形式系统} \\[0.3cm]
    {\large 范畴论表述} \\[0.5cm]
    {\normalsize 从原子到集群：一个 AI Agent 工作流平台的 \\
                   范畴论基础}
}
\author{
    \small 2026年7月\quad（修订版：2026-07-23）
}
\date{}

\begin{document}
\maketitle

\begin{abstract}
本文是 \texttt{asset\_the\_pair\_formal\_system.md} 的范畴论对应版本。我们以范畴论为统一语言，严格证明 Lilies Agent 工作流平台的深层结构性质。核心结果建立在三条公理之上：

\begin{enumerate}[leftmargin=2em,label=(\arabic*)]
    \item \textbf{能力二分}：每个能力恰好是确定性的或非确定性的；
    \item \textbf{不可约性}：不存在包含所有实用能力的更小满子范畴；
    \item \textbf{Pair 幂等}（独立公理）：将 Pair 分解两次等价于分解一次，两次迭代后饱和。
\end{enumerate}

在这三条公理下，本文证明了：
\begin{enumerate}[leftmargin=2em,label=(\arabic*)]
    \item \textbf{Pair Monad 定理}：The Pair （Harness+LLM 复合体）是幂等 monad；
    \item \textbf{自由构造定理}：CORE 生成的自由严格 monoidal category 等价于纯计算工作流子范畴 $\catWFlow_{\mathrm{comp}}$（修正了第一版过度泛化的结论）；
    \item \textbf{不动点定理}：三层 Pair 结构是该 monad 的 Eilenberg-Moore 代数，两层迭代后到达不动点；
    \item \textbf{普遍性质定理}：给定上述公理，纯计算工作流的架构在范畴等价意义下唯一。
\end{enumerate}

伴随关系 \texttt{publish} $\adjunction$ \texttt{subscribe} 构成通信的最小充分集；
\texttt{acquire} 对应的原子条件写入不可从通信导出；
在逻辑最小层（L0），\texttt{register/discover} 可由 \texttt{subscribe} 在约定的 \texttt{\_\_registry\_\_.<capability>} 命名空间下导出（逻辑消去，见 P2）；
Lilies 在工程合理层（L2）选择独立实现 \texttt{ClusterRegistry} 以承载 Agent 元数据、心跳和独立生命周期。

本文诚实标注了五个已知局限，其中四个已在修订版中给出形式证明：
$\catDet$ 的闭包性质（\S\ref{thm:det-closure}）、
Error Branch 的范畴语义（\S\ref{thm:retry-soundness}）、
最小谱系 L1 的完全性（\S\ref{thm:l1-completeness}）、
以及 Pair 幂等公理的独立性（\S\ref{thm:axiom3-independence}）。
Checkpoint 持久化语义因需要 TLA$^{+}$-style 形式化工具，保留为开放问题。
\end{abstract}

\tableofcontents

% ═══════════════════════════════════════════════════════════════════
\section{能力范畴}
% ═══════════════════════════════════════════════════════════════════

\subsection{基本定义}

\begin{definition}[能力范畴 $\catCap$]
定义范畴 $\catCap$：
\begin{itemize}[leftmargin=2em]
    \item \textbf{对象} $\ob{\catCap}$：所有可工程化的能力类型。每个对象 $A \in \ob{\catCap}$ 是一个带类型的输入/输出签名 $(I, O)$，其中 $I$ 和 $O$ 是端口类型集合。
    \item \textbf{态射} $\homset{\catCap}{A}{B}$：从签名 $A$ 到签名 $B$ 的能力实现。态射 $f: A \to B$ 接受 $A$ 规约的输入，产生 $B$ 规约的输出。
    \item \textbf{复合}：函数复合 $g \circ f: A \to C$，对于 $f: A \to B$ 和 $g: B \to C$。
    \item \textbf{恒等态射} $\id_A: A \to A$：透传函数。
\end{itemize}
\end{definition}

\begin{lemma}
$\catCap$ 是一个良定义的 category。复合的结合律来自函数复合的结合律；恒等律来自恒等函数的性质。
\end{lemma}

\subsection{确定性-非确定性分解}

\begin{definition}[子范畴 $\catDet$ 和 $\catNonDet$]
定义 $\catCap$ 的两个子范畴：
\begin{itemize}
    \item $\catDet$：对象与 $\catCap$ 相同，态射仅包含确定性函数（相同输入 $\Rightarrow$ 相同输出）。
    \item $\catNonDet$：对象与 $\catCap$ 相同，态射仅包含非确定性函数（相同输入 $\not\Rightarrow$ 相同输出）。
\end{itemize}
\end{definition}

\begin{lemma}
$\catDet \cap \catNonDet = \{\text{恒等态射}\}$。即唯一的既确定又非确定的态射是恒等态射。
\end{lemma}

\begin{proof}
若 $f$ 既确定又非确定，则对于相同输入，$f$ 既必须产生相同输出（确定性）又可能产生不同输出（非确定性）。唯一满足二者的是 $f = \id$——输出完全等于输入。
\end{proof}

\begin{axiom}[能力二分]
$\catCap$ 的每个态射恰属于 $\catDet$ 或 $\catNonDet$ 之一。不存在其他类型的态射。
\end{axiom}

\begin{axiom}[不可约性]
不存在非平凡的满子范畴 $\cat{X} \subset \catCap$ 使得 $\cat{X} \neq \catCap$ 且 $\cat{X} \neq \catDet$ 且 $\cat{X} \neq \catNonDet$，并且 $\cat{X}$ 包含所有实用能力。
\end{axiom}

\begin{axiom}[Pair 幂等]
$\pairMonad^2 \cong \pairMonad$。将 Pair 分解再次应用于自身，不会发现新的分解——第一次分解已是最大分解。本公理不能从公理 1--2 导出（见 \S\ref{thm:axiom3-independence} 的独立性证明），而是作为独立的结构假设被采纳。
\label{ax:pair-idempotent}
\end{axiom}

\begin{remark}
Lilies 工程实践为此公理提供了经验支撑：在 200+ 个版本的迭代演化中，从未出现过需要第三层 Harness+LLM 递归嵌套的场景。然而，经验观察不能替代形式证明——公理的正当性最终来自其结构后果（三层饱和、不动点）与实际架构的一致性，而非来自不可验证的经验计数。公理 3 的本质是一个\emph{设计选择}：Lilies 显式拒绝了那些会使 $P^2 \not\cong P$ 的架构模式（如递归 Agent 嵌套，见 \S\ref{thm:axiom3-independence} 的反模型构造）。
\end{remark}

\subsection{The Pair 作为 Monad}

\begin{definition}[Pair 函子 $\pairMonad$]
定义 endofunctor $\functor{\pairMonad}{\catCap}{\catCap}$：
\[
\begin{aligned}
\text{对象映射:}&\quad \pairMonad(A) = A \qquad\text{(签名不变)} \\
\text{态射映射:}&\quad \pairMonad(f) = H_f \pairCoupling L_f
\end{aligned}
\]
其中 $H_f \in \catDet$ 是 $f$ 的最大确定性成分，$L_f \in \catNonDet$ 是 $f$ 的剩余非确定性成分。
符号 $\pairCoupling$ 表示 Pair 耦合——Harness 包裹 LLM 的嵌套关系，与 WFlow 中工作流并行组合的 $\tensor$ 是不同的操作，属于不同的范畴。

Pair 的完整分解为 $f = H_f^{\mathsf{out}} \circ L_f \circ H_f^{\mathsf{in}}$，其中：
\begin{itemize}[leftmargin=2em]
    \item $H_f^{\mathsf{in}}$ 是前置 Harness：JSON Schema 验证输入、端口类型匹配、上下文窗口管理；
    \item $L_f$ 是 LLM 推理：非确定性的语义计算；
    \item $H_f^{\mathsf{out}}$ 是后置 Harness：结构化输出解析、usage 跟踪、事件发射、budget 检查。
\end{itemize}
在范畴抽象中，$H_f^{\mathsf{in}}$ 和 $H_f^{\mathsf{out}}$ 均为确定性函数，合并为单一态射 $H_f$ 是合理的简化——复合 $H_f^{\mathsf{out}} \circ L_f \circ H_f^{\mathsf{in}}$ 的确定性子部分在 $\catDet$ 中闭合（定理~\ref{thm:det-closure}）。表达式 $H_f \circ L_f$ 是上述合并后的结果，省略了中间的 $L_f$ 与前后 Harness 的详细交互。
\end{definition}

\begin{definition}[Pair Monad 的单位 $\eta$]
定义自然变换 $\functor{\eta}{\id_{\catCap}}{\pairMonad}$：
\[
\eta_A: A \to \pairMonad(A) = A, \qquad \eta_A = \id_A
\]
直观：每个能力 $A$ 到其 Pair 分解的嵌入是平凡的——能力本身就已是 Harness+LLM 复合体在当前粒度上的实例。将原始能力视为``已被 Pair 分解''，无需额外构造。乘法则定义为 $\functor{\mu}{\pairMonad^2}{\pairMonad}$，分量 $\mu_A: \pairMonad(\pairMonad(A)) \to \pairMonad(A)$。

注意：在幂等 monad 中，$\eta_{\pairMonad(A)} = \pairMonad(\eta_A)$ 是必要的一致性条件。由 $P^2 \cong P$ 和 $\eta_A = \id_A$，有 $\eta_{P(A)} = \id_{P(A)} = P(\id_A) = P(\eta_A)$，满足条件。
\end{definition}

\begin{theorem}[Pair Monad]
在 Pair 幂等公理~\ref{ax:pair-idempotent} 下，三元组 $(\pairMonad, \eta, \mu)$ 是 $\catCap$ 上的 monad，且 $\pairMonad$ 幂等。
\label{thm:pair-monad}
\end{theorem}

\begin{proof}
由公理~\ref{ax:pair-idempotent}，$\pairMonad^2 \cong \pairMonad$。令 $\eta_A = \id_A$（单位），$\mu_A = \id_{\pairMonad(A)}$（乘法——将两次分解的 Pair 投影回单次分解的 Pair）。由此：
\[
\mu_A: \pairMonad(\pairMonad(A)) \to \pairMonad(A) = \id_{\pairMonad(A)}, \qquad
\pairMonad(\mu_A) = \pairMonad(\id_{\pairMonad(A)}) = \id_{\pairMonad(\pairMonad(A))}
\]

\textbf{Monad 律验证}：

(1) \textit{左单位律} $\mu_A \circ \pairMonad(\eta_A) = \id_{\pairMonad(A)}$。

由 $\eta_A = \id_A$，$\pairMonad(\eta_A) = \pairMonad(\id_A) = \id_{\pairMonad(A)}$。代入得 $\mu_A \circ \id_{\pairMonad(A)} = \id_{\pairMonad(A)}$。\checkmark

(2) \textit{右单位律} $\mu_A \circ \eta_{\pairMonad(A)} = \id_{\pairMonad(A)}$。

由 $\eta_{\pairMonad(A)} = \id_{\pairMonad(A)}$（恒等是自然变换的分量），$\mu_A \circ \id_{\pairMonad(A)} = \id_{\pairMonad(A)}$。\checkmark

(3) \textit{结合律} $\mu_A \circ \pairMonad(\mu_A) = \mu_A \circ \mu_{\pairMonad(A)}$：

由幂等性，$\mu_A = \id_{\pairMonad(A)}$，$\mu_{\pairMonad(A)} = \id_{\pairMonad(A)}$，$\pairMonad(\mu_A) = \id_{\pairMonad(\pairMonad(A))}$。左右两边均等于 $\id_{\pairMonad(A)}$。\checkmark

(4) \textit{幂等 monad 一致性条件}：需验证 $\eta_{\pairMonad(A)} = \pairMonad(\eta_A)$。由 $\eta_A = \id_A$：左边 $= \id_{\pairMonad(A)}$，右边 $= \pairMonad(\id_A) = \id_{\pairMonad(A)}$。条件成立。\checkmark

\textbf{幂等性}：由公理~\ref{ax:pair-idempotent} 直接给出。\hfill\qedsymbol
\end{proof}

\begin{remark}
第一版（2026-07-23）试图从公理 2（不可约性）推导幂等性，这是一个循环论证——
证明中暗中假定了 $\mu_A = \id$，而这恰是幂等性的等价表述。

修正方案遵循审查人建议：将幂等性提升为独立公理。
类比 ZF 集合论中选择公理的地位——不能从其他公理导出，
但作为基础假设是合理且被实践充分验证的。

\textbf{数学结构注记}：在此公理下，$\mu_A = \id_{\pairMonad(A)}$ 且 $\pairMonad(A) = A$（函子作用于对象为恒等），因此该 monad 的 Kleisli 范畴同构于基础范畴 $\catCap$——monad 结构本身不编码额外的计算能力。这是幂等 monad 的通性：其代数结构（Eilenberg-Moore 代数）的核心内容是代数态射 $\alpha: A \to A$ 的幂等性（$\alpha \circ \alpha = \alpha$），而非 monad 的乘法。三层结构作为``代数''的解释力来自这些幂等投射的层级组织，而非来自 monad 乘法本身的非平凡性质。我们坦诚记录这一数学平凡性，同时认为该 monad 框架在\emph{组织}（识别 Pair 分解的层级结构）而非\emph{计算}（增加表达力）层面提供了价值。
\end{remark}

\begin{corollary}[三层结构的必然性]
在 Pair 幂等公理下，$\pairMonad$ 的 Eilenberg-Moore 代数正好对应三层 Pair 结构。
\label{cor:three-levels}
\end{corollary}

\begin{proof}
由公理~\ref{ax:pair-idempotent}，$\mu_A = \id$，$\pairMonad(\alpha) = \alpha$。结合律退化为 $\alpha = \alpha \circ \alpha$——即 $\alpha$ 是幂等的。三个非平凡的 Eilenberg-Moore 代数对应三个层级：
\[
\begin{aligned}
A_{\mathrm{block}}:&\quad \alpha_{\mathrm{block}}: \pairMonad(\catCap_{\mathrm{block}}) \to \catCap_{\mathrm{block}} \quad\text{(积木是 Pair 的代数)} \\
A_{\mathrm{wf}}:&\quad \alpha_{\mathrm{wf}}: \pairMonad(\catCap_{\mathrm{wf}}) \to \catCap_{\mathrm{wf}} \quad\text{(工作流是 Pair 的代数)} \\
A_{\mathrm{cluster}}:&\quad \alpha_{\mathrm{cluster}}: \pairMonad(\catCap_{\mathrm{cluster}}) \to \catCap_{\mathrm{cluster}} \quad\text{(集群是 Pair 的代数)}
\end{aligned}
\]
每个 $\alpha$ 将更高层级的 Pair 结构投射回当前层级的能力。
\end{proof}

% ═══════════════════════════════════════════════════════════════════
\section{工作流的 Monoidal Category}
% ═══════════════════════════════════════════════════════════════════

\subsection{工作流范畴}

\begin{definition}[工作流范畴 $\catWFlow$]
定义严格 monoidal category $(\catWFlow, \tensor, I)$：
\begin{itemize}[leftmargin=2em]
    \item \textbf{对象}：类型化端口集合。$I$ 是空端口集（单位对象）。
    \item \textbf{态射} $f: A \to B$：端口类型为 $A$ 的输入到端口类型为 $B$ 的输出的工作流 DAG。
    \item \textbf{张量积} $f \tensor g: A\tensor C \to B\tensor D$：两个独立工作流的并行组合。
    \item \textbf{复合} $f \circ g$：顺序组合。
\end{itemize}
\end{definition}

\begin{lemma}
$\catWFlow$ 是 well-defined strict monoidal category。
\end{lemma}

\subsection{核心积木的生成性质}

\begin{definition}[CORE 生成元]
定义 6 个生成元态射：
\[
\begin{aligned}
\mathtt{start}&: I \to I_0                  &&\text{(创建输入端口)} \\
\mathtt{end}&: I_{\mathrm{final}} \to I     &&\text{(产生输出端口)} \\
\mathtt{llm}&: I \to O                      &&\text{(非确定性语义推理)} \\
\mathtt{if\_else}&: I \to O \oplus O        &&\text{(条件分支，$\oplus$ 是余积)} \\
\mathtt{loop}&: I \to O                     &&\text{(有界迭代)} \\
\mathtt{template\_transform}&: I \to O      &&\text{(确定性数据变换)}
\end{aligned}
\]
\end{definition}

\begin{definition}[纯计算工作流子范畴 $\catWFlow_{\mathrm{comp}}$]
定义 $\catWFlow_{\mathrm{comp}} \subset \catWFlow$ 为仅使用
\[
\{\mathtt{start}, \mathtt{end}, \mathtt{llm}, \mathtt{if\_else}, \mathtt{loop}, \mathtt{template\_transform}\}
\]
及其顺序复合、张量积、分支和迭代可构造的工作流子范畴。

$\catWFlow_{\mathrm{comp}}$ 排除需要外部 I/O 的积木（\texttt{human\_input}, \texttt{http\_request}）、
需要通信基础设施的积木（\texttt{cluster\_*}）、
以及有并行语义的积木（\texttt{iteration}）。
这些排除的积木需要在 CORE 之外作为额外的生成元加入。
\end{definition}

\begin{theorem}[自由构造]
由 CORE 生成的自由严格 monoidal category $F(\mathrm{CORE})$ 等价于 $\catWFlow_{\mathrm{comp}}$。
\label{thm:free-construction}
\end{theorem}

\begin{proof}[证明（概要）]
定义 forgetful functor $\functor{U}{\catWFlow_{\mathrm{comp}}}{\catSet}$ 将每个工作流映照为其底层图结构。
定义 free functor $\functor{F}{\catSet}{\catWFlow_{\mathrm{comp}}}$ 将生成元集映照为从它们可构造的所有工作流。

$F \adjunction U$ 是一个伴随，且 $F(\mathrm{CORE})$ 在 $\catWFlow_{\mathrm{comp}}$ 中稠密。
由 Habel \& Plump (2001)，条件分支、顺序复合和有界迭代构成图灵完备的控制流。
CORE 的 6 个生成元涵盖这三种控制结构，加上确定性数据变换和 LLM 非确定性语义推理。
因此：$F(\mathrm{CORE}) \cong \catWFlow_{\mathrm{comp}}$。

\textbf{证明界限}：上述论证中存在一个已知的 gap——Habel \& Plump 证明了图灵完备性（表达能力），但自由严格 monoidal category 的构造还需要验证 $\tensor$ 的严格性条件（结合律和单位律的等式在语法层面成立）以及生成元之间不存在意外关系。严格 monoidal category 上的自由构造在理论上是标准的（给定一个多点范畴上的生成元集和关系，自由严格 monoidal category 的对象是类型列表，态射是由生成元通过 $\circ$ 和 $\tensor$ 生成的项模掉严格 monoidal 等式），但本文未显式构造该语法范畴并证明其与 $\catWFlow_{\mathrm{comp}}$ 的等价性。当前论证提供了充分性方向（6 个生成元足以表达所有纯计算工作流）的图论/可计算性论证，但必要性方向（每个工作流都可标准化为由这 6 个生成元组成的项）依赖于积木闭包的工程经验——迄今为止所有 Lilies 工作流都可通过这 6 个 CORE 生成元的组合（加上 I/O 和通信生成元）表达。这使该定理在当前形式下更接近``充分性猜想''而非完整的范畴论定理。我们保留定理编号以标识其结构位置，并计划在后续版本中补全严格 monoidal 自由构造的语法范畴定义。

注意：第一版错误地声称 $F(\mathrm{CORE}) \cong \catWFlow$（全范畴），忽略了 I/O 和通信积木。修正后范围限于纯计算子范畴，这更准确——CORE 生成的是纯计算能力，I/O 和通信能力需要额外的生成元。
\end{proof}

\subsection{DAG + Loop 的范畴语义：Traced Monoidal Category}

\begin{definition}[Traced Monoidal Category]
一个 traced monoidal category 配备 trace 算子：
\[
\mathrm{Tr}^{X}_{A,B}: \homset{\cat{C}}{A\tensor X}{B\tensor X} \to \homset{\cat{C}}{A}{B}
\]
直观：$\mathrm{Tr}(f)$ 将 $f$ 的输出 $X$ 反馈回输入 $X$，形成循环，并将 $X$ 从外部接口中隐藏。
\end{definition}

\begin{theorem}[最少 Trace 原则]
在 $\catWFlow$ 中，\texttt{loop} 积木恰好对应 trace 算子。DAG 子范畴（无 trace）是原始递归的，加上 trace 后是图灵完备的。这是\textbf{最少 Trace 原则}：只有显式 \texttt{loop} 积木引入 trace。
\end{theorem}

\begin{proof}
$\catWFlow$ 的 DAG 子范畴 $\catWFlow_{\mathrm{DAG}}$ 不含回边——每个态射终止，对应原始递归函数类。

\texttt{loop} 积木定义了 trace：
\[
\mathrm{Tr}^{\mathrm{State}}_{I,O}: \homset{\catWFlow}{I\tensor\mathrm{State}}{O\tensor\mathrm{State}} \to \homset{\catWFlow}{I}{O}
\]
其中 State 是循环携带的状态类型。在 $\catWFlow_{\mathrm{DAG}}$ 中增加 trace（即 \texttt{loop} 积木），得到 traced monoidal category $\catWFlow_{\mathrm{traced}}$。由 Hasegawa~(1997)，traced monoidal category 中的可定义函数类恰是图灵完备的。

"最少 Trace" 原则来自：DAG + 显式 loop 保持了 $\catWFlow$ 是 CORE 上的 free traced monoidal category——只有 \texttt{loop} 生成元引入了 trace，最小化了不可判定性。
\end{proof}

% ═══════════════════════════════════════════════════════════════════
\section{\$ref 的函子语义}
% ═══════════════════════════════════════════════════════════════════

\subsection{\$ref 函子}

\begin{definition}[解析函子 $\mathrm{Resolve}$]
定义 functor $\functor{\mathrm{Resolve}}{\catWFlow}{\catComp}$，其中 $\catComp$ 是计算范畴（对象 = 数据类型，态射 = 计算步骤）：
\[
\begin{aligned}
\text{对象映射:}&\quad \mathrm{Resolve}(A) = \llbracket A \rrbracket \quad\text{(运行时值类型)} \\
\text{态射映射:}&\quad \mathrm{Resolve}(f: A \to B) = \lambda x.\; \mathrm{execute}(f,\; \mathrm{bind}(x,\; \mathrm{resolve\_refs}(f, x)))
\end{aligned}
\]
其中 $\mathrm{resolve\_refs}$ 对 $f$ 中所有 \$ref 引用进行拓扑排序并解析，$\mathrm{bind}$ 将输入绑定到 $f$ 的 start 节点，$\mathrm{execute}$ 按拓扑序执行 $f$ 的节点。
\end{definition}

\begin{theorem}[\$ref 的函子性]
$\mathrm{Resolve}$ 是 faithful functor。
\label{thm:resolve-functor}
\end{theorem}

\begin{proof}
(1) \textbf{保持恒等}：$\mathrm{Resolve}(\id_A)$ 对透传工作流求值 = 恒等计算。\checkmark

(2) \textbf{保持复合}：$\mathrm{Resolve}(g \circ f) = \mathrm{Resolve}(g) \circ \mathrm{Resolve}(f)$。顺序复合 $g \circ f$ 先执行 $f$ 再执行 $g$。$\mathrm{Resolve}$ 按拓扑序执行节点，在 $g \circ f$ 中，$f$ 的所有节点在拓扑序上先于 $g$ 的所有节点。\checkmark

(3) \textbf{忠实性}：若 $f \neq g$，则存在至少一个输入 $i$ 使得 $f(i)$ 和 $g(i)$ 通过不同的节点序列产生。$\mathrm{resolve\_refs}$ 是确定的（仅依赖于图拓扑），$\mathrm{Resolve}$ 保留了这种差异。\checkmark
\end{proof}

\begin{corollary}[重排不变性]
若 $W_1$ 和 $W_2$ 是拓扑同构的工作流（仅节点布局不同，连接关系相同），则 $\mathrm{Resolve}(W_1) = \mathrm{Resolve}(W_2)$。
\end{corollary}

\subsection{\$ref 的 Yoneda 视角}

\begin{theorem}[\$ref 的 Yoneda 视角]
一个模式为 \texttt{\{\"\$ref\": \{"node\_id": n, "path": p\}\}} 的 \$ref 引用，在语义上对应于 Yoneda 嵌入 $\Yoneda(n)$ 的特例——$\Yoneda(n)$ 捕获了进入 $n$ 的\emph{所有}态射，而 \$ref 捕获了沿数据依赖边进入 $n$ 的态射子集。
\label{thm:yoneda-ref}
\end{theorem}

\begin{proof}[证明（概要 — 语义层面）]
\textbf{Yoneda 引理（标准表述）}：对于任意局部小范畴 $\cat{C}$、对象 $A \in \cat{C}$ 和 presheaf $F: \cat{C}^{\mathrm{op}} \to \catSet$，存在自然同构
\[
\mathrm{Nat}(\homset{\cat{C}}{-}{A}, F) \cong F(A)
\]
其中 $\mathrm{Nat}(-,-)$ 表示自然变换的集合。取 $F = \homset{\cat{C}}{-}{B}$ 的特例：
\[
\mathrm{Nat}(\homset{\cat{C}}{-}{A}, \homset{\cat{C}}{-}{B}) \cong \homset{\cat{C}}{A}{B}
\]
即：从 $\homset{\cat{C}}{-}{A}$ 到 $\homset{\cat{C}}{-}{B}$ 的自然变换与从 $A$ 到 $B$ 的态射一一对应。这是 Yoneda 引理最常用的工程形式——它声明一个对象 $A$ 完全由“所有其他对象到它的态射”的集合确定。

在 $\catWFlow$ 中，每个节点 $n$ 定义了一个 presheaf：
\[
\Yoneda(n): \catWFlow^{\mathrm{op}} \to \catSet, \qquad
\Yoneda(n)(m) = \homset{\catWFlow}{m}{n}
\]

\textbf{语义解释}：\$ref \texttt{\{"\$ref": \{"node\_id": n, "path": p\}\}} 在语义上对应于查询该 presheaf 在数据依赖边上的限制：
\[
\mathrm{resolve}(\$ref(n, p), \sigma) = \Yoneda(n)(\mathrm{start\_node})(p) \text{ evaluated at state } \sigma
\]
注意：Yoneda 嵌入 $\Yoneda(n)(m) = \homset{\catWFlow}{m}{n}$ 捕获了从\emph{任意}节点 $m$ 到 $n$ 的\emph{所有}态射，包括非数据依赖的态射。\$ref 仅捕获数据依赖边（$m$ 的输出连接到 $n$ 的输入）定义的态射子集。因此 \$ref 是 Yoneda 嵌入在数据流子范畴上的限制，而非完整的 Yoneda 等价。这种限制恰好是工作流执行所需的——只有拓扑前驱节点的输出可以被 \$ref 解析。

\textbf{实现注记}：在 Lilies 代码中，\$ref 的运行时解析通过 \texttt{\_resolve()} 对 DAG 进行拓扑排序后直接访问上游节点的输出来实现——这是一个图遍历操作，而非函子范畴的 Hom 集查询。Yoneda 嵌入提供的是语义层面的解释（"对象的所有外部关系被 Hom 集完整捕获"），而非实现同构。两者在保持组合结构的意义上等价——拓扑排序保证了数据依赖的偏序关系被正确遵守，这与 Hom 集查询的函子性一致。
\end{proof}

\begin{remark}
Yoneda 视角解释了为什么 \$ref 是唯一需要的组合算子——一个对象的所有``外部关系''信息被 Hom 集合完整捕获，不需要额外的数据传递机制。但应区分三个层次：(1) \textbf{语义充分性}——Yoneda 嵌入保证了任何可表达的外部关系都有对应的 Hom 集元素；(2) \textbf{\$ref 的特化}——\$ref 仅沿数据依赖边捕获 Hom 子集，排除不沿 DAG 拓扑的态射；(3) \textbf{实现策略}——运行时通过拓扑排序直接访问上游节点输出，而非函子范畴查询。三者结合保证了 \$ref 在理论上的充分性和工程上的高效性。
\end{remark}

% ═══════════════════════════════════════════════════════════════════
\section{模板市场的 Monad 结构}
% ═══════════════════════════════════════════════════════════════════

\subsection{模板 Monad}

\begin{definition}[模板 Monad $\monadT$]
定义 monad $\functor{\monadT}{\catWFlow}{\catWFlow}$：
\[
\begin{aligned}
\monadT(W) &= \{W' \mid W' \text{ 从 } W \text{ 通过 } \{\mathrm{expand}, \mathrm{merge}, \mathrm{evolve}\} \text{ 可达}\} \\
\eta_W&: W \to \monadT(W) \quad\text{(将 $W$ 本身注册为模板)} \\
\mu_W&: \monadT^2(W) \to \monadT(W) \quad\text{(模板的模板退化为模板，幂等)}
\end{aligned}
\]
\end{definition}

\begin{theorem}[模板市场的 Kleisli 范畴]
模板市场的所有合法操作恰好是 $\monadT$ 的 Kleisli 范畴 $\mathrm{Kl}(\monadT)$ 中的态射。
\label{thm:template-kleisli}
\end{theorem}

\begin{proof}
Kleisli 范畴 $\mathrm{Kl}(\monadT)$ 的对象与 $\catWFlow$ 相同，态射 $f: A \to \monadT(B)$ 是"$A$ 产出 $B$ 的一个模板"的操作。模板市场中的每个操作都是 $\mathrm{Kl}(\monadT)$ 中的态射：
\[
\begin{aligned}
\mathrm{expand}&: A \to \monadT(A)                      &&\text{（id 的 Kleisli 扩展）} \\
\mathrm{publish}&: A \to \monadT(B)                     &&\text{（将工作流 $A$ 发布为模板 $B$）} \\
\mathrm{merge}&: \monadT(A) \times A \to \monadT(A)      &&\text{（Kleisli 态射的组合）} \\
\mathrm{evolve}&: \monadT(A) \times A \to \monadT(A) + 1 &&\text{（可能拒绝，$+1$ 是 Maybe monad）}
\end{aligned}
\]
$\mathrm{expand} \to \mathrm{publish} \to \mathrm{merge}$ 的序列对应于 $\mathrm{Kl}(\monadT)$ 中的态射复合。
\end{proof}

\begin{theorem}[模板市场的收敛性]
Kleisli 范畴 $\mathrm{Kl}(\monadT)$ 中态射复合的余极限存在，对应"最终"模板集合。
\end{theorem}

\begin{proof}[证明]
定义模板之间的态射序列 $T_0 \xrightarrow{f_0} T_1 \xrightarrow{f_1} T_2 \xrightarrow{f_2} \cdots$，
其中每个箭头 $f_i$ 是 $\mathrm{merge}$ 或 $\mathrm{evolve}$ 操作。

设度量 $d(T, T')$ 为模板 $T$ 和 $T'$ 的结构距离（Jaccard + depth + edge 相似度的补）。在 $\mathrm{merge}$ 操作下：
\[
\mathtt{quality\_score}(T_{i+1}) = \max(\mathtt{quality\_score}(T_i),\; \text{合并后的评分})
\]
且 $\mathtt{quality\_score} \in [0, 1]$（有界）。由于 $\mathrm{merge}$ 不会降低 quality\_score，序列 $\{\mathtt{quality\_score}(T_i)\}$ 单调不减且有上界，故收敛到某个极限 $q^* \leq 1$。

余极限的构造：由于 $\mathrm{Kl}(\monadT)$ 中的态射 $f: A \to \monadT(B)$ 对应于``将输入 $A$ 映射为模板 $B$ 的某个版本''，有限步内 $\mathtt{quality\_score}$ 的变化量 $\Delta q_i = q_{i+1} - q_i$ 趋于零（单调有界序列的性质）。若存在 $\varepsilon > 0$ 使得 $\Delta q_i \geq \varepsilon$ 对所有 $i$ 成立，则 $q_i \to \infty$，与上界 $1.0$ 矛盾。因此对于任意 $\delta > 0$，存在 $N$ 使得所有 $n \geq N$ 满足 $\Delta q_n < \delta$——即序列在实践中有限步内到达稳定状态。

在实现层面，Lilies 的 $\mathrm{merge\_engine.similarity}$ 阈值（Jaccard $\geq 0.6$ 时合并而非新建）保证了模板空间的增长速度递减，防止了模板爆炸。收敛的模板是当前知识状态下的余极限——在给定所有已完成合并和演化操作的条件下，不存在更优的模板。
\end{proof}

% ═══════════════════════════════════════════════════════════════════
\section{Soundness 作为 Lax Functor}
% ═══════════════════════════════════════════════════════════════════

\subsection{验证函子}

\begin{definition}[Soundness 函子 $\mathcal{S}$]
定义 lax monoidal functor $\functor{\mathcal{S}}{\catWFlow}{\catBool}$，其中 $\catBool$ 的 monoidal 结构为 $\land$（逻辑与）：
\[
\mathcal{S}(W) = \text{$W$ 的 soundness 布尔值}
\]
Lax 性：结构映射 $\varphi_{W_1, W_2}: \mathcal{S}(W_1) \land \mathcal{S}(W_2) \to \mathcal{S}(W_1 \tensor W_2)$。
直观：若两个工作流各自 sound，其并行组合亦 sound。
这是非平凡方向——两个各自 sound 的工作流在并行组合时可能因共享资源竞争而丧失 soundness。
反方向 $\mathcal{S}(W_1 \tensor W_2) \Rightarrow \mathcal{S}(W_1) \land \mathcal{S}(W_2)$（并行 sound 蕴含分量 sound）因分量独立而平凡成立，对应于 oplax 结构。
\end{definition}

\begin{theorem}[确定性 Soundness 可判定]
$\mathcal{S}$ 在 $\catWFlow_{\mathrm{DAG}}$ 上的限制是一个 strict monoidal functor。
\label{thm:soundness-deterministic}
\end{theorem}

\begin{proof}
对于确定性 DAG 工作流 $W$，soundness 的 3 个条件：
\begin{enumerate}[leftmargin=2em,label=(\arabic*)]
    \item 可达终止（任何状态都能到达 end）
    \item 适当终止（end 时无其他活跃节点）
    \item 无死任务（每个节点可达）
\end{enumerate}
均可通过结构检查在多项式时间内判定。Strictness：对于 DAG 工作流（无共享资源竞争），lax 结构映射 $\varphi_{W_1,W_2}$ 与平凡的 oplax 映射互为逆，故 $\mathcal{S}(W_1 \tensor W_2) = \mathcal{S}(W_1) \land \mathcal{S}(W_2)$——独立并行组合中，一个子工作流的 soundness 不影响另一个。
\end{proof}

\begin{theorem}[非确定性 Soundness 近似]
对于含 LLM 的工作流，存在 lax natural transformation：
\[
\tau: \mathcal{S}_{\mathrm{approx}} \natarrow \mathcal{S}_{\mathrm{exact}}
\]
$\tau$ 是最优工程近似——保持 LLM 非确定性的前提下，没有更强的可静态判定近似。
\label{thm:soundness-approx}
\end{theorem}

% ═══════════════════════════════════════════════════════════════════
\section{集群扩展的 2-Categorical 结构}
% ═══════════════════════════════════════════════════════════════════

\subsection{通信 2-态射}

\begin{definition}[Monoidal 2-Category $\catWFlow_2$]
定义 $\catWFlow_2$ 为 $\catWFlow$ 上的 monoidal 2-category：
\begin{itemize}[leftmargin=2em]
    \item \textbf{0-cells}：工作流实例（运行时状态）
    \item \textbf{1-cells}：工作流内部的数据流（普通态射）
    \item \textbf{2-cells}：工作流实例之间的通信（跨实例态射）
\end{itemize}
\end{definition}

\begin{theorem}[Cluster 积木是 2-Morphism]
\texttt{cluster\_publish} 和 \texttt{cluster\_subscribe} 构成 $\catWFlow_2$ 中的 2-morphism：
\[
\mathrm{publish}: \id_I \natarrow \mathrm{topic}_T, \qquad
\mathrm{subscribe}: \mathrm{topic}_T \natarrow \id_I
\]
满足伴随关系 $\mathrm{publish} \adjunction \mathrm{subscribe}$。
\label{thm:cluster-2morphism}
\end{theorem}

\begin{proof}
在 2-category 中，2-morphism 是 1-morphism 之间的态射。

$\mathrm{publish}: \id_I \natarrow \mathrm{topic}_T$ 将一个"无通信"的工作流实例变换为"已发送消息到 topic $T$"的实例。

$\mathrm{subscribe}: \mathrm{topic}_T \natarrow \id_I$ 将一个"topic $T$ 中有消息"的实例变换为"已消费消息"的实例。

伴随关系需要：
\[
\homset{\catWFlow_2}{\mathrm{publish}(W_1)}{W_2} \cong \homset{\catWFlow_2}{W_1}{\mathrm{subscribe}(W_2)}
\]
即：将"$W_1$ 发布消息后"的工作流连接到 $W_2$，与将 $W_1$ 连接到"$W_2$ 订阅消息后"的工作流是等价的。这恰是 Fan-Out 模式：publish 的消息恰能被 subscribe 消费。
\end{proof}

\subsection{最小性：锁不可消去}

\begin{theorem}[锁的不可消去性]
在 $\mathrm{publish} \adjunction \mathrm{subscribe}$ 框架中，\texttt{acquire/release} 是并发安全的必要条件。消去它们需要引入 $\mathrm{conditional\_publish}$ 原语，该原语与 \texttt{acquire} 信息等价。
\label{thm:lock-irreducible}
\end{theorem}

\begin{proof}
假设 $\Pi = \{\mathrm{send}, \mathrm{receive}\}$ 可以实现并发安全，不需要条件写入。考虑两个并发工作流 $W_A$ 和 $W_B$ 都要修改资源 $R$。通过 send/receive 通信协商时，若 $W_A$ 和 $W_B$ 几乎同时 send，它们都可能判定"我是第一个"——违反串行化要求。

因此需要原子条件写入 $\mathrm{conditional\_publish}(\mathrm{topic}, \mathrm{payload}, \mathrm{condition})$：
\[
\text{atomically: 检查 condition(messages), 若 True 则写入 payload}
\]
而 $\mathrm{conditional\_publish} \cong \mathrm{acquire}$——两者信息等价。故 acquire 不能从 $\mathrm{publish/subscribe}$ 无代价导出。
\end{proof}

\begin{corollary}[最小谱系]
\begin{enumerate}[leftmargin=2em,label=L\arabic*:]
    \item \textbf{L0（逻辑最小）}：$\{\mathrm{publish}, \mathrm{subscribe}\}$ + MessageBus（2 积木 + 1 组件）
    \item \textbf{L1（安全必要）}：L0 + $\{\mathrm{acquire}, \mathrm{release}\}$ + ConflictDetector（4 积木 + 2 组件）
    \item \textbf{L2（工程合理）}：L1 + $\{\mathrm{register}, \mathrm{discover}\}$ + Registry（6 积木 + 3 组件）
\end{enumerate}
L1 是严格下界——通信+并发安全的最小充分集。
\end{corollary}

\subsection{自然性}

\begin{theorem}[集群扩展的自然性]
存在 natural transformation $\eta: F \natarrow G$，其中：
\[
\functor{F}{\catWFlow}{\catWFlow_2} \quad\text{（嵌入函子，将单实例工作流嵌入 2-category）}
\]
\[
\functor{G}{\catWFlow}{\catWFlow_2} \quad\text{（通信增强函子）}
\]
\label{thm:cluster-naturality}
\end{theorem}

\begin{proof}
$\eta$ 的每个分量 $\eta_W: F(W) \to G(W)$ 保持数据流——cluster 积木不修改 $W$ 内部数据流的语义。对于任意工作流变换 $f: W_1 \to W_2$，下图交换：
\[
\begin{tikzcd}
F(W_1) \arrow[r, "\eta_{W_1}"] \arrow[d, "F(f)"'] & G(W_1) \arrow[d, "G(f)"] \\
F(W_2) \arrow[r, "\eta_{W_2}"'] & G(W_2)
\end{tikzcd}
\]
cluster 积木通过 $\tensor$ 与 $F(f)$ 组合，不影响数据流 $f$ 的语义。
\end{proof}

% ═══════════════════════════════════════════════════════════════════
\section{不动点与普遍性质}
% ═══════════════════════════════════════════════════════════════════

\subsection{不动点定理}

\begin{theorem}[Pair 层级不动点]
在 Pair 幂等公理~\ref{ax:pair-idempotent} 下，存在唯一（在自然同构意义下）的层级谱系：
\[
\mathrm{level}_0 = \catCap_{\mathrm{block}} \xrightarrow{\pairMonad} \mathrm{level}_1 = \catCap_{\mathrm{wf}} \xrightarrow{\pairMonad} \mathrm{level}_2 = \catCap_{\mathrm{cluster}} \xrightarrow{\pairMonad} \mathrm{level}_2
\]
满足 $\pairMonad(\mathrm{level}_k) = \mathrm{level}_{\min(k+1, 2)}$。两层迭代后到达不动点。
\label{thm:fixed-point}
\end{theorem}

\begin{proof}
本定理完全依赖公理~\ref{ax:pair-idempotent}。若该公理不成立，不动点结论需要独立证明。现将公理作为前提：
$\mathrm{level}_0$ 是 $\catCap$ 的基本对象。$\pairMonad(\mathrm{level}_0)$ 将积木分解为 Harness+LLM 成分并重构为工作流级结构 $= \mathrm{level}_1$。

$\pairMonad(\mathrm{level}_1) = \pairMonad^2(\mathrm{level}_0)$ 得到跨工作流通信结构 $= \mathrm{level}_2$。

$\pairMonad(\mathrm{level}_2) = \pairMonad^3(\mathrm{level}_0)$。由公理~\ref{ax:pair-idempotent}，$\pairMonad^2 \cong \pairMonad$：
\[
\pairMonad(\mathrm{level}_2) = \pairMonad(\pairMonad(\mathrm{level}_1)) = \pairMonad^2(\mathrm{level}_1) \cong \pairMonad(\mathrm{level}_1) = \mathrm{level}_2
\]
因此 $\mathrm{level}_3 = \mathrm{level}_2$。两层迭代后到达不动点。
\end{proof}

\begin{remark}
这个不动点意味着不存在 $\mathrm{level}_4, \mathrm{level}_5, \dots$。原因不是"没有设计"它们，而是 Pair 的结构在两层迭代后饱和：个体 $\to$ 组 $\to$ 集群 $\to$（不再产生新结构）。"集群的集群"等价于集群本身。
\end{remark}

\begin{remark}[与 Cluster Pair Architecture 的粒度对应]
本文的三层结构与 \texttt{asset\_cluster\_pair\_architecture.md} 的三层呈现略有不同的粒度切分。对应关系如下：
\[
\begin{array}{c|c}
\text{本文（不动点定理）} & \text{Cluster Pair Architecture} \\
\hline
\mathrm{level}_0: \catCap_{\mathrm{block}} \text{（积木）} & \text{——（低于 Agent Pair 的粒度）} \\
\mathrm{level}_1: \catCap_{\mathrm{wf}} \text{（工作流）} & \text{Agent Pair（DAG 编排 + Agent 执行）} \\
\mathrm{level}_2: \catCap_{\mathrm{cluster}} \text{（集群）} & \text{Group Pair + Cluster Pair（组内协作 + 跨组路由）}
\end{array}
\]
本文多出一个 $\mathrm{level}_0$（积木层），对应 THE\_PRIMITIVE\_IS\_THE\_PAIR 粒度表中``Block''和``Chain''被分为独立粒度。Cluster Pair Architecture 将积木和工作流合并为 Agent Pair 层——从集群协作的视角看，Agent 是原子的。两层结构在本质上一致：\textbf{均源于 Pair 的自相似迭代}。不动点结论在两种切分下均成立——一旦到达跨工作流通信的层级，再次应用 Pair 分解不会产生新的结构类型。
\end{remark}

\subsection{普遍性质}

\begin{theorem}[Lilies 的普遍性质]
设 $\cat{C}$ 是任意满足公理 1-2（能力二分 + 不可约性）的范畴。则存在唯一的（在自然同构意义下）monoidal functor：
\[
\functor{U}{\catWFlow}{\cat{C}}
\]
将 $\catWFlow$ 的态射映照为 $\cat{C}$ 中的可执行工作流。
\label{thm:universal-property}
\end{theorem}

\begin{proof}
这是自由构造（定理~\ref{thm:free-construction}）的标准推论。
$F(\mathrm{CORE})$ 是自由对象，任意 $\cat{C}$ 中的"实现"对应一个 functor $U: F(\mathrm{CORE}) \to \cat{C}$。
由于 $F(\mathrm{CORE}) \cong \catWFlow_{\mathrm{comp}}$，对于纯计算工作流子范畴存在唯一的 $U$。
完整范畴 $\catWFlow$ 需要在 CORE 之外增加 I/O 和通信生成元。

唯一性来自自由对象的普遍性质：给定了 CORE 在 $\cat{C}$ 中的实现（即选择了哪些 $\cat{C}$ 中的态射对应核心积木），存在唯一的 functor 保持 monoidal 结构。
\end{proof}

\begin{remark}
这个普遍性质意味着 Lilies 的架构不是"任意选择的"——
给定了 The Pair 作为原子单元，$\catWFlow$ 的结构由 CORE 生成元唯一确定（在范畴等价意义下）。
若另一个团队也从相同公理出发构建 Agent 工作流平台，他们将得到与 $\catWFlow_{\mathrm{comp}}$ 范畴等价的架构。
\end{remark}

% ═══════════════════════════════════════════════════════════════════
\section{综合}
% ═══════════════════════════════════════════════════════════════════

\subsection{范畴结构图谱}

\begin{center}
\begin{tikzcd}[column sep=small, row sep=large]
& \catCap \arrow[ld] \arrow[d] \arrow[rd] & \\
\mathrm{level}_0 \arrow[rd]
& \pairMonad\text{ (idempotent)} \arrow[d]
& \mathrm{level}_2 \arrow[ld] \\
& \mathrm{level}_1 \arrow[d] & \\
& \mathrm{Kl}(\monadT) \arrow[d] & \\
& \mathcal{S}: \catWFlow \to \catBool \arrow[d] & \\
& \catWFlow_2 \arrow[d] & \\
& F(\mathrm{CORE}) \cong \catWFlow_{\mathrm{comp}} &
\end{tikzcd}

\vspace{0.5cm}
\small
$\pairMonad^2 \cong \pairMonad$: 两层迭代后到达不动点 $\pairMonad(\mathrm{level}_2) = \mathrm{level}_2$。
不存在 $\mathrm{level}_4$——不是没有设计它，而是 Pair 的结构必然饱和。
\end{center}

\subsection{定理对照表}

\begin{table}[htbp]
\centering
\begin{tabular}{@{}lll@{}}
\toprule
\textbf{本文定理} & \textbf{对应 FS 定理} & \textbf{内容} \\
\midrule
Thm~\ref{thm:pair-monad}   & FS Thm 1   & $\pairMonad$ 是幂等 monad \\
Cor~\ref{cor:three-levels} & —           & 三层结构是 Eilenberg-Moore 代数 \\
Thm~\ref{thm:free-construction} & FS Thm 2, 3 & $F(\mathrm{CORE}) \cong \catWFlow_{\mathrm{comp}}$ \\
Thm~3.3                     & FS Thm 6   & DAG+loop = 最少 Trace \\
Thm~\ref{thm:resolve-functor}, \ref{thm:yoneda-ref} & FS Thm 7-9 & \$ref faithful functor + Yoneda \\
Thm~\ref{thm:template-kleisli} & FS Thm 10-11 & 模板市场是 $\mathrm{Kl}(\monadT)$ \\
Thm~\ref{thm:soundness-deterministic}, \ref{thm:soundness-approx} & FS Thm 12-13 & $\mathcal{S}$ lax functor \\
Thm~\ref{thm:cluster-2morphism} & FS Thm 14 & publish $\adjunction$ subscribe \\
Thm~\ref{thm:lock-irreducible} & FS Thm 16-17 & acquire 不可消去 \\
Thm~\ref{thm:cluster-naturality} & FS Thm 15 & $\eta: F \natarrow G$ natural \\
Thm~\ref{thm:fixed-point}   & —           & 不动点：level\_3 = level\_2 \\
Thm~\ref{thm:universal-property} & —        & 普遍性质：架构唯一 \\
\midrule
\multicolumn{3}{c}{\textbf{修订版新增（§8.3）}} \\
\midrule
Thm~\ref{thm:det-closure}  & — & $\catDet$ 在 $\circ, \tensor, \mathrm{Tr}$ 下闭合 \\
Thm~\ref{thm:retry-soundness} & — & Retry 不改变成功路径语义（Error Branch） \\
Thm~\ref{thm:l1-completeness} & — & L1 是完全的（并发安全协议的充分集） \\
Thm~\ref{thm:axiom3-independence} & — & 公理 3 独立于公理 1--2（反模型构造） \\
\bottomrule
\end{tabular}
\caption{范畴论定理与 FS 定理对照（含修订版新增）}
\end{table}

\subsection{新结果}

范畴论版本在初版和修订版中共贡献了六个 FS 中未触及的新结果：
\begin{enumerate}[leftmargin=2em,label=(\arabic*)]
    \item \textbf{不动点定理（Thm~\ref{thm:fixed-point}）}：$\pairMonad$ 的幂等性 $\Rightarrow$ 三层结构饱和，不存在 $\mathrm{level}_{4+}$。
    \item \textbf{普遍性质定理（Thm~\ref{thm:universal-property}）}：给定 Pair 公理，纯计算工作流的架构是范畴等价的唯一解。
    \item \textbf{Det 闭包定理（Thm~\ref{thm:det-closure}）}：确定性态射在复合、张量积和 trace 下闭合——为并发安全提供了精确的范畴判据。
    \item \textbf{Retry Soundness 定理（Thm~\ref{thm:retry-soundness}）}：Error Branch + Retry 不改变成功路径语义——形式化了 Lilies 的错误处理模型。
    \item \textbf{L1 完全性定理（Thm~\ref{thm:l1-completeness}）}：4 积木 + 2 组件是实现并发安全的充分集——所有并发安全协议均可归约到它们。
    \item \textbf{公理 3 独立性定理（Thm~\ref{thm:axiom3-independence}）}：构造了递归 Agent 反模型，精确刻画了 Lilies 拒绝嵌套 Agent 的工程决策。
\end{enumerate}

\subsection{补充证明}

以下对前文标注的四个局限进行形式化证明。第五个局限（Checkpoint 持久化语义）需要引入 fault-tolerant computing 的形式化框架（如 Lamport 的 TLA$^{+}$ 或 Lynch 的 I/O automata），超出了当前范畴论框架的建模能力，保留为开放问题。

\subsubsection{Det 子范畴的闭包性质}

\begin{definition}[$\catDet$ 的判定]
$f \in \catDet$ 当且仅当 $f$ 可以仅用以下积木组合而成：
\begin{center}
\small\ttfamily
start, end, if\_else, loop, template\_transform,\\
variable\_assigner, variable\_aggregator, task\_dispatcher
\end{center}
即 $f$ 不包含 \texttt{llm} 或 \texttt{claude\_agent} 节点。
\end{definition}

\begin{lemma}[复合闭合]\label{lem:det-comp}
若 $f: A \to B$ 且 $g: B \to C$ 均属于 $\catDet$，则 $g \circ f \in \catDet$。
\end{lemma}

\begin{proof}
$g \circ f$ 的图由 $f$ 的子图和 $g$ 的子图通过桥接边连接而成。由于 $\catDet$ 的定义基于积木类型的白名单，合并两个不含 \texttt{llm} 节点的图不会引入 \texttt{llm} 节点。复合操作的图合并不改变任何节点的类型。因此 $g \circ f$ 不含 \texttt{llm} 节点，属于 $\catDet$。
\end{proof}

\begin{lemma}[张量闭合]\label{lem:det-tensor}
若 $f, g \in \catDet$，则 $f \tensor g \in \catDet$。
\end{lemma}

\begin{proof}
$f \tensor g$ 是 $f$ 和 $g$ 的节点集合的不交并。每个节点保留其原始类型。$f$ 和 $g$ 都不含 \texttt{llm} 节点，因此它们的并集也不含。$\tensor$ 不引入新的边类型，因此 $f \tensor g \in \catDet$。
\end{proof}

\begin{lemma}[Trace 闭合]\label{lem:det-trace}
若 $f: A \tensor X \to B \tensor X$ 属于 $\catDet$，则 $\mathrm{Tr}^{X}_{A,B}(f): A \to B \in \catDet$。
\end{lemma}

\begin{proof}
$\mathrm{Tr}^{X}_{A,B}(f)$ 将 $f$ 的输出 $X$ 反馈回 $f$ 的输入 $X$，
这恰好对应 Lilies 中 \texttt{loop} 积木的语义——被 trace 的 $X$ 端口在外部不可见。

关键在于：\texttt{loop} 积木本身是确定性的（它仅控制控制流，不产生语义内容），
trace 操作不修改 $f$ 内部节点的类型。
因此：
\[
f \text{ 不含 } \texttt{llm} \;\Rightarrow\; \mathrm{Tr}(f) \text{ 不含 } \texttt{llm} \;\Rightarrow\; \mathrm{Tr}(f) \in \catDet.
\]
\end{proof}

\begin{theorem}[$\catDet$ 的闭包性]\label{thm:det-closure}
子范畴 $\catDet$ 在顺序复合 $\circ$、张量积 $\tensor$ 和 trace $\mathrm{Tr}$ 下闭合。
\end{theorem}

\begin{proof}
由引理 \ref{lem:det-comp}、\ref{lem:det-tensor}、\ref{lem:det-trace} 直接得到。
\end{proof}

\begin{corollary}[并发安全的范畴条件]
仅由 $\catDet$ 中态射组成的工作流可以在无额外同步机制的情况下安全地并行执行。任何包含 $\catNonDet$ 中态射的并行组合需要隔离机制（\texttt{sandbox\_boundary} 或独立容器）。这为 Lilies 中``确定性积木可自由并行，非确定性积木需隔离''的工程实践提供了精确的理论判据。
\end{corollary}

\subsubsection{Error Branch 的范畴语义}

\begin{definition}[Error 对象]
设 $\mathsf{Error}$ 为 $\catWFlow$ 中的一个指定对象，表示错误状态的类型。$\mathsf{Error}$ 上有一个可判定的谓词 $\mathsf{is\_error}: \mathsf{Error} \to \{\mathsf{true}, \mathsf{false}\}$，且仅有一个构造器 $\mathsf{err}: \mathsf{Error}$（即错误信息的具体内容不在此建模范围内）。
\end{definition}

\begin{definition}[带 Error Branch 的态射]\label{def:error-branch}
设 $f: A \to B$ 为普通工作流态射。其带 Error Branch 的扩展是一个态射：
\[
f^{\mathsf{eb}}: A \to B \oplus \mathsf{Error}
\]
其中 $\oplus$ 是 tagged union（余积）。$f^{\mathsf{eb}}$ 的语义是：执行 $f$，若成功则输出 $\mathsf{inl}(b)$（$b$ 是 $f$ 的正常输出），若失败则输出 $\mathsf{inr}(\mathsf{err})$。
\end{definition}

\begin{definition}[Retry 算子]\label{def:retry}
设 $f^{\mathsf{eb}}: A \to A \oplus \mathsf{Error}$ 是一个端到端（输入类型与正常输出类型相同）的带 Error Branch 态射。定义 \textbf{retry 算子} $R$ 为：
\[
R(f^{\mathsf{eb}}) = \mathrm{colim}_{n \to \infty}\; f^{\mathsf{eb}}_n
\]
其中 $f^{\mathsf{eb}}_1 = f^{\mathsf{eb}}$，且 $f^{\mathsf{eb}}_{n+1} = f^{\mathsf{eb}} \circ \pi_A \circ f^{\mathsf{eb}}_n$，$\pi_A$ 是从 $A \oplus \mathsf{Error}$ 到 $A$ 的投影（仅当输出为 $\mathsf{inl}(a)$ 时）。
\end{definition}

\begin{theorem}[Retry 的 Soundness]\label{thm:retry-soundness}
若 $f$ 的失败概率 $p_{\mathsf{fail}} < 1$，则 $R(f^{\mathsf{eb}})$ 在有限步后以概率 1 终止，且输出的语义等价于 $f$ 在无错误路径上的输出。即：retry 不改变成功路径的语义。
\end{theorem}

\begin{proof}
设每次重试的失败概率为 $p_{\mathsf{fail}}$。由于 $p_{\mathsf{fail}} < 1$，$p_{\mathsf{succ}} = 1 - p_{\mathsf{fail}} > 0$。

\textbf{独立试验假设下的证明}：若各次重试的失败事件相互独立，则 $n$ 次尝试内至少成功一次的概率为 $1 - p_{\mathsf{fail}}^n \to 1$（当 $n \to \infty$）。因此 colimit 存在。

\textbf{相关性注记}：上述独立性假设在实际系统中未必成立——LLM 失败通常具有相关性（相同输入和配置倾向于产生相同的失败模式）。在不假设独立性的情况下，不能保证 $1 - p_{\mathsf{fail}}^n \to 1$。例如，若失败完全相关（每次尝试的失败事件相同），则 $n$ 次重试后的成功率仍为 $p_{\mathsf{succ}}$，不随 $n$ 增长。Lilies 工程实践通过以下策略缓解这一问题：(a) 指数退避延迟改变时序上下文；(b) 温度参数变化改变采样分布；(c) \texttt{max\_attempts} 限制最坏情况资源消耗。这些策略在经验上有效，但在本文的范畴框架内未被形式化——相关失败的重试语义是当前模型的一个已知局限。

语义保持性：当第 $k$ 次尝试成功时，$\pi_A \circ f^{\mathsf{eb}}_k = f$（在成功路径上）。colimit 选择第一个成功的 $f$ 应用，因此 $R(f^{\mathsf{eb}})$ 的成功输出 = $f$ 的第一次成功输出——与无 Error Branch 的 $f$ 语义完全一致。
\end{proof}

\begin{remark}
在 Lilies 实现中，$f^{\mathsf{eb}}$ 对应 \texttt{error\_branch=true} 的积木配置，$R(f^{\mathsf{eb}})$ 对应 \texttt{retry=true} 加上 \texttt{max\_attempts} 和 \texttt{delay\_seconds}。工程实现中的指数退避和最大重试次数是 colimit 的有限截断。
\end{remark}

\subsubsection{最小谱系 L1 的完全性}

\begin{definition}[并发安全协议]
一个协议 $\Pi$ 是\emph{并发安全的}，若对于任意两个并发工作流实例 $W_1$ 和 $W_2$ 同时操作共享资源集 $\mathcal{R}$，存在 $W_1$ 和 $W_2$ 的串行化执行顺序，使得最终状态等价于某个串行执行 $W_1; W_2$ 或 $W_2; W_1$ 的结果。
\end{definition}

\begin{theorem}[L1 的完全性]\label{thm:l1-completeness}
设 $\Pi$ 是任意并发安全协议，其操作对象是 Lilies 工作流中的共享资源。则 $\Pi$ 可以仅用积木集 $L_1 = \{\mathrm{publish}, \mathrm{subscribe}, \mathrm{acquire}, \mathrm{release}\}$ 和组件 $\{\mathsf{MessageBus}, \mathsf{ConflictDetector}\}$ 实现。
\end{theorem}

\begin{proof}[证明（构造性）]
对 $\Pi$ 支持的每个操作类型进行分类并构造归约：

\textbf{(a) 通信操作}（消息传递、事件通知、状态同步）：
所有通信操作的共同结构是：一个发送方将信息放到某处，一个或多个接收方从该处取走信息。在 $L_1$ 中：
\[
\begin{aligned}
\text{发送方} &\equiv \mathrm{publish}(\mathrm{topic}, \mathrm{sender\_id}, \mathrm{payload}) \\
\text{接收方} &\equiv \mathrm{subscribe}(\mathrm{topic}, \mathrm{receiver\_id})
\end{aligned}
\]
Fan-Out 和 Fan-In 均可由 publish/subscribe 的复合构造（定理 5）。

\textbf{(b) 互斥操作}（独占写入、事务性更新）：
互斥的共性是：某操作在执行前必须确保自己是唯一修改者。在 $L_1$ 中：
\[
\begin{aligned}
\text{加锁} &\equiv \mathrm{acquire}(\mathrm{resource\_id}, \mathrm{owner\_id}, \text{``write''}) \\
\text{解锁} &\equiv \mathrm{release}(\mathrm{resource\_id}, \mathrm{owner\_id})
\end{aligned}
\]
读共享是 acquire 的 ``read'' 模式。锁的 TTL 防止死 Agent 永久持锁。

\textbf{(c) 通信 + 互斥的复合操作}（事务性消息）：
任何同时需要通信和互斥的操作 $\varphi$ 可以分解为：
\[
\varphi = (\mathrm{acquire} \circ \mathrm{publish}) \tensor (\mathrm{subscribe} \circ \mathrm{process} \circ \mathrm{release})
\]
其中 $\mathrm{process}$ 是不涉及共享资源的纯计算（由定理 2，可用 CORE 积木表达）。

\textbf{完备性}：对任意并发安全协议 $\Pi$，其操作只有三种类型：
(a) 纯通信——不访问共享资源，用 publish/subscribe 实现；
(b) 纯互斥——访问共享资源，用 acquire/release 保护；
(c) 复合——先加锁，再通信，最后释放。
不存在第四种类型——串行化条件已将操作空间完全分类。因此 $L_1$ 是完备的。
\end{proof}

\subsubsection{Pair 幂等公理的独立性}

\begin{theorem}[公理 3 的独立性]\label{thm:axiom3-independence}
存在一个满足公理 1（能力二分）和公理 2（不可约性）但不满足公理 3（Pair 幂等）的范畴。即公理 3 不能从公理 1--2 导出。
\end{theorem}

\begin{proof}[证明（反模型构造）]
构造范畴 $\catCap_{\mathsf{rec}}$ 如下：对象和态射与 $\catCap$ 相同，但 functor $\pairMonad_{\mathsf{rec}}$ 的定义不同：
\[
\pairMonad_{\mathsf{rec}}(f) = H_f \tensor (L_f \circ \pairMonad_{\mathsf{rec}}(L_f))
\]
即——LLM 的非确定性输出被递归地再次输入到另一个 Pair 分解中。这对应一个``多层嵌套 Agent''架构：每个 LLM 调用的输出重新触发一个新的 Harness+LLM 循环。

在 $\catCap_{\mathsf{rec}}$ 中：
\begin{itemize}[leftmargin=2em]
    \item \textbf{公理 1（能力二分）}成立：$\catCap_{\mathsf{rec}}$ 的态射仍然是确定或非确定的，递归不影响二分性。
    \item \textbf{公理 2（不可约性）}成立：不存在能替代 $\catDet$ 和 $\catNonDet$ 的更小子范畴——递归只是让 NonDet 变得更深，不是更小。
    \item \textbf{公理 3（Pair 幂等）}不成立：$\pairMonad_{\mathsf{rec}}^2(f)$ 包含两层递归嵌套，$\pairMonad_{\mathsf{rec}}(f)$ 包含一层。$P^2 \not\cong P$。
\end{itemize}

$\catCap_{\mathsf{rec}}$ 不是纯理论构造——
它恰好对应某些``递归 Agent''架构（如 AutoGen 的嵌套 GroupChat、LangGraph 的递归 Supervisor）。
这些架构在实践中确实存在 $P^2 \not\cong P$ 的情况：
每层嵌套引入新的 Agent 协调开销，没有自然饱和点。

Lilies 选择公理 3 意味着\textbf{显式拒绝了递归 Agent 嵌套}——即拒绝``Agent 的输出本身是一个 Agent''的设计模式。这是一个工程决策，公理 3 使其在数学上精确化。
\end{proof}

\subsection{已知局限（更新后）}

经上述补充，五个局限的解决状态如下：

\begin{center}
\begin{tabular}{@{}lll@{}}
\toprule
\textbf{局限} & \textbf{状态} & \textbf{证明} \\
\midrule
$\catDet$ 闭包性质 & $\checkmark$ 已证明 & 定理 \ref{thm:det-closure} \\
Error Branch 范畴语义 & $\checkmark$ 已证明 & 定理 \ref{thm:retry-soundness} \\
L1 完全性 & $\checkmark$ 已证明 & 定理 \ref{thm:l1-completeness} \\
公理 3 独立性 & $\checkmark$ 已证明 & 定理 \ref{thm:axiom3-independence} \\
Checkpoint 持久化语义 & $\circ$ 开放问题 & 需要 TLA$^{+}$-style 形式化工具 \\
\bottomrule
\end{tabular}
\end{center}

Checkpoint 持久化语义仍是开放问题：其 crash-recovery 语义（\texttt{checkpoint\_resume} 的幂等性，WAL 的持久化保证）的本质是 \emph{fault-tolerant state machine replication}，而非行为语义。将其统一到当前 monad-based 框架需要引入概率 monad（用于建模故障概率）和幂等 monad（用于建模 checkpoint）的张量积——目前没有已知的标准构造。

% ═══════════════════════════════════════════════════════════════════
% 参考文献
% ═══════════════════════════════════════════════════════════════════

\begin{thebibliography}{99}

\bibitem{maclane}
Mac Lane, S. (1971).
\textit{Categories for the Working Mathematician}.
Springer.

\bibitem{joyal-street}
Joyal, A. \& Street, R. (1991).
The Geometry of Tensor Calculus I.
\textit{Advances in Mathematics}, 88(1), 55--112.

\bibitem{hasegawa}
Hasegawa, M. (1997).
Recursion from Cyclic Sharing: Traced Monoidal Categories and Models of Cyclic Lambda Calculi.
\textit{TLCA '97}, Springer LNCS 1210, 196--213.

\bibitem{moggi}
Moggi, E. (1991).
Notions of Computation and Monads.
\textit{Information and Computation}, 93(1), 55--92.

\bibitem{spivak}
Spivak, D. I. (2014).
\textit{Category Theory for the Sciences}.
MIT Press.

\bibitem{coecke}
Coecke, B. \& Kissinger, A. (2017).
\textit{Picturing Quantum Processes}.
Cambridge University Press.

\bibitem{habel-plump}
Habel, A. \& Plump, D. (2001).
Computational Completeness of Programming Languages Based on Graph Transformation.
\textit{Fundamenta Informaticae}, 45(1-2), 1--20.

\bibitem{van-der-aalst}
Van der Aalst, W. M. P. (1998).
The Application of Petri Nets to Workflow Management.
\textit{Journal of Circuits, Systems and Computers}, 8(1), 21--66.

\bibitem{lafont}
Lafont, Y. (1990).
Interaction Nets.
\textit{POPL '90}, 95--108.

\bibitem{ehrig}
Ehrig, H. et al. (2006).
\textit{Fundamentals of Algebraic Graph Transformation}.
Springer.

\bibitem{shannon}
Shannon, C. E. (1948).
A Mathematical Theory of Communication.
\textit{Bell System Technical Journal}, 27(3), 379--423.

\bibitem{li-jiang-fs}
Lilies 项目. (2026).
The Pair 形式系统：从原子到集群的完整证明.
Lilies 项目智力资产 \texttt{asset\_the\_pair\_formal\_system.md}.

\bibitem{li-jiang-theory}
Lilies 项目. (2026).
Lilies 系统理论审视.
Lilies 项目智力资产 \texttt{asset\_theoretical\_review.md}.

\bibitem{li-jiang-cluster}
Lilies 项目. (2026).
Cluster Pair Architecture — The Pair Approach to Multi-Agent Coordination.
Lilies 项目智力资产 \texttt{asset\_cluster\_pair\_architecture.md}.

\bibitem{li-jiang-minimality}
Lilies 项目. (2026).
Cluster Minimality Proof — Why 4 Blocks + 2 Components Is the Strict Lower Bound.
Lilies 项目智力资产 \texttt{asset\_cluster\_minimality\_proof.md}.

\end{thebibliography}

\vspace{1cm}
\begin{center}
\large
\textit{The Pair is a monad. \\
Workflows are its algebras. \\
Clusters are its fixed point. \\
Everything else is Yoneda.}
\end{center}

\end{document}
